\documentclass[english, parskip=full]{scrartcl}
\usepackage[english]{babel}  % Sprache auf Deutsch 
\usepackage[utf8]{inputenc}  % Kodierung der Datei
\usepackage[left=2cm, right=2cm, top=2cm, bottom=3cm, footskip=1.5cm]{geometry}
\usepackage{color}
\usepackage{graphicx, gensymb}
\usepackage{amsfonts, cancel, amssymb}
\usepackage{amsmath, amsthm, array, esint}
\usepackage{bm} % bold math symbols, e.g. \bm{\sigma} etc.
\usepackage{booktabs, multirow}  % schönere Tabellen
%\usepackage{scrlayer-scrpage}    % Manipulation des Seitenstils
%\pagestyle{scrheadings}  % Stil für die Seite setzen
\usepackage{tabularx}
\usepackage{setspace}
%\automark{section}  % Abschnittsnamen als Seitenbeschriftung 
%\setheadsepline{.5pt}    % Kopzeile durch Linie abtrennen
\usepackage[hidelinks]{hyperref}  % Links und weitere PDF-Features
\usepackage{typearea, tikz, pgffor, verbatim}
\usetikzlibrary{calc, arrows.meta,arrows, graphs, quotes, positioning, angles, babel}
\usepackage[cache=false, outputdir=build]{minted}
\usemintedstyle{vs}
\usepackage{placeins} % provides \FloatBarrier

\definecolor{bg}{rgb}{0.9,0.9,0.9}
\newcommand\numberthis{\addtocounter{equation}{1}\tag{\theequation}}
\usepackage{svg}
\usepackage{siunitx}
\usepackage{physics}
\usepackage{tensor}
\usepackage{slashed}
\usepackage{bbm}
\usepackage[compat=1.1.0]{tikz-feynman}

\usepackage{caption}
\captionsetup{
    % width=.8\textwidth,
    % indention=1cm,
    margin=0.5cm,
    format=plain,
    labelfont=bf
}
%\usepackage{subcaption}
%\subcaptionsetup{
%    indention=0cm,
%    format=hang,
%    margin=0.5cm,
%    labelfont=normalfont
%}
\usepackage[english,capitalise]{cleveref}
\usepackage[
    backend=biber,
    sorting=none,
    style=phys,
    giveninits=true,
    sortcites=true,
    citestyle=numeric-comp,
    % url=false,
    % doi=false,
    biblabel=brackets,
    % eprint=false,
    maxcitenames=3]{biblatex}
\usepackage{csquotes}
%\addbibresource{bibliography.bib}

\usepackage{mathtools}
% use \abs for absolute values
% use \abs for \left ... \right version and \abs[\bigg for \bigg version etc.
\let\abs\undefined
\let\bra\undefined
\let\braket\undefined
\DeclarePairedDelimiter\abs{\vert}{\vert}
\DeclarePairedDelimiter\kla{(}{)}
\DeclarePairedDelimiter\klb{[}{]}
\DeclarePairedDelimiter\klc{\{}{\}}
\DeclarePairedDelimiter\bra{\langle}{\rangle}
\DeclarePairedDelimiterX\brb[2]{\langle #1}{\rangle}{\delimsize\vert #2 }
\DeclarePairedDelimiterX\brc[3]{\langle #1}{\rangle}{\delimsize\vert #2 \delimsize\vert #3 }

\renewcommand{\i}{\text{i}}
\newcommand{\e}{\text{e}}
\DeclareMathOperator{\sgn}{sgn}
\newcommand{\kom}[1]{\overset{\substack{#1 \\ |}}{=}}
\newcommand{\koml}[2]{\overset{\substack{#1 \\ |}}{#2}}
\newcommand{\intx}[4]{\int_{#1}^{#2} #3 \,\mathrm{d}#4}
\newcommand{\intr}[2]{\int_{\mathbb{R}} #1 \, \mathrm{d}#2}
\newcommand{\intrnd}[3]{\int_{\mathbb{R}^{#1}} #2 \, \mathrm{d}^{#1}#3}
\newcommand{\intrpi}[2]{\int_{\mathbb{R}} #1 \, \frac{\mathrm{d}#2}{2\pi}}
\newcommand{\intrndpi}[3]{\int_{\mathbb{R}^{#1}} #2 \, \frac{\mathrm{d}^{#1}#3}{(2\pi)^{#1}}}
\newcommand{\oper}[2]{\vert #1 \rangle \langle #2 \vert}
\newcommand{\Text}[1]{\text{#1}}    % this is relevant because \text does not autocomplete to the default '\text{@}' but rather '\textbf' etc.
\newcommand{\powe}[1]{\e^{#1}}
\newcommand{\pow}[2]{{#1}^{#2}}
\newcommand{\Ket}[1]{\vert #1 \rangle}
\newcommand{\Bra}[1]{\langle #1 \vert}
\newcommand*{\Fourier}{\textsc{Fourier}\ }
\newcommand*{\F}{\mathcal{F}}
\DeclareMathOperator{\arsinh}{arsinh}
\DeclareMathOperator{\arcosh}{arcosh}
\DeclareMathOperator{\artanh}{artanh}
\DeclareMathOperator{\sinc}{sinc}
\DeclareMathOperator{\cscc}{cscc}
\DeclareMathOperator{\sinhc}{sinhc}
\DeclareMathOperator{\cschc}{cschc}
\newcommand{\conv}{\mathop{\scalebox{3.5}{\raisebox{-0.38ex}{\makebox[0.5\width][c]{$\ast$}}}}\limits}
\newcommand*{\unity}{\text{\usefont{U}{bbold}{m}{n}1}}   % operator one from :: https://tex.stackexchange.com/questions/566780/import-mathbb1-character-from-the-unicode-math-package

\renewcommand{\d}{\text{d}}
\renewcommand{\r}{\text{r}}
\renewcommand{\t}{\text{t}}
\renewcommand{\a}{\text{a}}
\renewcommand{\o}{\text{o}}
\renewcommand{\F}{\text{F}}
\renewcommand{\c}{\text{c}}
\newcommand{\A}{\text{A}}
\newcommand{\B}{\text{B}}
\newcommand{\R}{\text{R}}
\newcommand{\M}{\text{M}}
\newcommand{\m}{\text{m}}
\newcommand{\s}{\text{s}}
\newcommand{\f}{\text{f}}
\newcommand{\K}{\text{K}}
\newcommand{\hc}{\text{h.c.}}
\newcommand{\kf}{k_\F}
\newcommand{\vf}{v_\F}
\newcommand{\const}{\text{const.}}
\renewcommand{\L}{\text{L}}
\newcommand{\gray}[1]{\bgroup\color{white!40!black}#1\egroup}
\newcommand{\TODO}[1]{\bgroup\color{red!60!black}#1\egroup}
\newcommand\QnA[1]{\bgroup\color{blue}#1\egroup}
\newcommand\highlight[1]{\bgroup\color{green!60!black}#1\egroup}

\newcommand{\cabx}[3]{\begin{smallmatrix}\gray{A}&\gray{:}&#1 \\ \gray{B}&\gray{:}&#2 \\ \gray{\Xi}&\gray{:}&#3\end{smallmatrix}}
\newcommand{\zabx}[3]{\begin{smallmatrix}\gray{A}&\gray{:}&#1 \\ \gray{B}&\gray{:}&#2 \\ \gray{Z}&\gray{:}&#3\end{smallmatrix}}
\newcommand{\abx}[3]{\begin{smallmatrix}#1 \\ #2 \\ #3\end{smallmatrix}}

\newcommand{\normord}[1]{
  {:\mathrel{\mspace{1mu}#1\mspace{1mu}}:}
}
\newcommand{\varnormord}[1]{
  {\mathrel{\mspace{1mu}\substack{\ast\\[-1.5pt] \ast}#1\substack{\ast\\[-1.5pt] \ast}\mspace{1mu}}}
}

% punctuation styles (see https://tex.stackexchange.com/questions/56063/how-to-properly-typeset-footnotes-superscripts-after-punctuation-marks)
\let\origfootnote\footnote
\renewcommand{\footnote}[1]{\kern.06em\origfootnote{#1}}
\newcommand{\punctfootnote}[1]{\kern-.06em\origfootnote{#1}}

% Multiline equations
\newcommand{\bsplit}[1]{\begin{aligned}[t]#1\end{aligned}}

\newcommand*{\titel}{Linear two-atomic chain}
\newcommand*{\autor}{Joachim Schwardt}

\hypersetup{pdfauthor={\autor}, pdftitle={\titel}}

%\titlehead{titlehead}
%\subject{SUBJECT}
\title{\titel}
\author{\autor}
\date{\today}

\begin{document}

%\begin{titlepage}
\maketitle  % Titel setzen
%\tableofcontents  % Inhaltsverzeichnis setzen
%\end{titlepage}

\section{Bosonization}
As before the model is described by $H=H_0+H_\Text{int}$ where
\begin{align}
H_0 &= \sum_{\kappa=\A,\B} \sum_k \vf k \varnormord{c_{\kappa,k}^\dagger c_{-\kappa,k}},\\
H_\Text{int} &= \sum_{\kappa=\A,\B} \sum_{j=0}^{N-1} U_\kappa n_{\kappa,j} n_{\kappa,j+1}.
\end{align}
The free part can be diagonalized by $c_{\kappa} = \frac{1}{\sqrt{2}} \kla*{c_{\R} + \kappa c_{\L}}$.
This basis rotation yields
\begin{align}
H_0 &= \sum_{\ell=\R,\L} \sum_k \vf k \frac{1}{2} \varnormord{\kla*{c_{\R,k}^\dagger + \kappa c_{\L,k}^\dagger} \kla*{c_{\R,k} - \kappa c_{\L,k}}} = \sum_{\ell=\R,\L} \sum_k \ell\vf k\varnormord{c_{\ell,k}^\dagger c_{\ell,k}}.
\end{align}
The density operator becomes $n_{\kappa} = \frac{1}{2} \kla*{c_{\R}^\dagger + \kappa c_{\L}^\dagger} \kla*{c_{\R} + \kappa c_{\L}} = \frac{1}{2} \sum_\ell \kla*{n_{\ell} + \kappa c_{\ell}^\dagger c_{-\ell}}$.
Inserting into the interaction yields
\begin{align}
H_\Text{int} &= \frac{1}{4} \sum_{\kappa=\A,\B} \sum_{j=0}^{N-1} U_\kappa \sum_{\ell,\ell'} \kla*{n_{\ell,j} + \kappa c_{\ell,j}^\dagger c_{-\ell,j}} \kla*{n_{\ell',j+1} + \kappa c_{\ell',j+1}^\dagger c_{-\ell',j+1}} = H_+ + H_-
\end{align}
where introducing $U_\pm = \frac{1}{2} \kla*{U_\A \pm U_\B}$ gives
\begin{align}
H_+ &= \frac{U_+}{2} \sum_{\ell,\ell'} \sum_{j=0}^{N-1} \kla*{n_{\ell,j} n_{\ell',j+1} + c_{\ell,j}^\dagger c_{-\ell,j} c_{\ell',j+1}^\dagger c_{-\ell',j+1}},\\
H_- &= \frac{U_-}{2} \sum_{\ell,\ell'} \sum_{j=0}^{N-1} \kla*{n_{\ell,j} c_{\ell',j+1}^\dagger c_{-\ell',j+1} + c_{\ell,j}^\dagger c_{-\ell,j} n_{\ell',j+1}}.
\end{align}
Going to a continuum theory via $c_j\rightarrow \sqrt{a}\psi(x)$ where $a$ is the lattice spacing and gives
\begin{align}
H_\pm &= \frac{U_\pm a}{2} \intr{\sum_{\ell,\ell'} \kla*{\mathcal{H}_{\pm,\ell,\ell'}^\Text{I} + \mathcal{H}_{\pm,\ell,\ell'}^\Text{II}}}{x}
\end{align}
where we defined the Hamilton densities
\begin{align}
\mathcal{H}_{+,\ell,\ell'}^\Text{I} &= \psi_{\ell}^\dagger(x) \psi_{\ell}(x) \psi_{\ell'}^\dagger(x+a) \psi_{\ell'}(x+a),\\
\mathcal{H}_{+,\ell,\ell'}^\Text{II} &= \psi_{\ell}^\dagger(x) \psi_{-\ell}(x) \psi_{\ell'}^\dagger(x+a) \psi_{-\ell'}(x+a),\\
\mathcal{H}_{-,\ell,\ell'}^\Text{I} &= \psi_{\ell}^\dagger(x) \psi_{\ell}(x) \psi_{\ell'}^\dagger(x+a) \psi_{-\ell'}(x+a),\\
\mathcal{H}_{-,\ell,\ell'}^\Text{II} &= \psi_{\ell}^\dagger(x) \psi_{-\ell}(x) \psi_{\ell'}^\dagger(x+a) \psi_{\ell'}(x+a).
\end{align}
The Fourier components relate to the real-space operators via the transform
\begin{align}
c_k &= \frac{1}{\sqrt{L}} \intr{\powe{-\i kx} \psi(x)}{x}
\end{align}
which, in the thermodynamic limit $L\rightarrow\infty$ and using partial integration, allows us to rewrite the Hamiltonian as
\begin{align}
H_0 &= \sum_{\ell,k} \ell \vf k  \frac{1}{L}\intr{\intr{\powe{\i kx} \powe{-\i kx'} \varnormord{\psi_\ell^\dagger(x) \psi_\ell(x')}}{x'} }{x} = \sum_\ell \ell\vf \intr{ \varnormord{\psi_\ell^\dagger(x) (-\i\partial_x) \psi_\ell(x)} }{x}.
\end{align}
The bosonization identity we will use is $\psi_\ell(x) = \frac{\eta_\ell}{\sqrt{L}} \normord{\powe{-\i\varphi_\ell(x)}}$.
Here $\eta_\ell$ obey the Majorana algebra and will for convenience be represented by $\eta_\R = \sigma_x$ and $\eta_\L = \sigma_z$.
The bosonic fields may be split into a part containing creation and a part containing annihilation operators via $\varphi_\ell = \varphi_\ell^{(+)} + \varphi_\ell^{(-)}$.
These satisfy the commutation relations
\begin{align}
\klb*{\varphi_\ell^{(+)}(x), \varphi_{\ell'}^{(-)}(x')} &= \lim_{\alpha\rightarrow 0^+} \delta_{\ell,\ell'} \log\klb*{\frac{2\pi}{L} \kla*{\alpha + \i\ell(x-x')}}.
\end{align}
\subsection{Free Hamiltonian}
To bosonize the free Hamiltonian we consider the point-splitting
\begin{align}
\psi_\ell^\dagger(x) \psi_\ell(x+\epsilon) &= \frac{1}{L} \normord{\powe{\i\varphi_\ell(x)}} \normord{\powe{-\i\varphi_\ell(x+\epsilon)}} = \frac{1}{L} \normord{\powe{-\i\varphi_\ell(x+\epsilon) + \i\varphi_\ell(x)}} \powe{\klb*{\i\varphi_\ell^{(-)}(x), -\i\varphi_\ell^{(+)}(x+\epsilon)}} \\
&= \frac{1}{L} \normord{\powe{-\i\sum_{n\ge 1} \frac{1}{n!} \epsilon^n \partial_x^n\varphi_\ell(x)}} \frac{L}{2\pi\i\ell \epsilon} \\
&= \frac{-\i\ell}{2\pi\epsilon} \normord{\klb*{1 -\i \epsilon \partial_x\varphi_\ell - \i \frac{\epsilon^2}{2} \partial_x^2\varphi_\ell - \frac{\epsilon^2}{2} \kla*{\partial_x\varphi_\ell}^2 + \dots} }.
\end{align}
Removing the divergent constant by subtracting the ground-state expectation value then yields
\begin{align}
\lim_{\epsilon\rightarrow 0} \varnormord{\psi_\ell^\dagger(x) \psi_\ell(x+\epsilon)} &= \frac{-\ell}{2\pi} \partial_x\varphi_\ell(x),\\
\lim_{\epsilon\rightarrow 0} \varnormord{\psi_\ell^\dagger(x)\partial_\epsilon \psi_\ell(x+\epsilon)} &= \frac{\i\ell}{4\pi} \klb*{\i \partial_x^2\varphi_\ell + \normord{\kla*{\partial_x\varphi_\ell}^2} }.
\end{align}
Inserting into the Hamiltonian and assuming that the fields vanish at infinity we obtain the well-known expression for a clean Luttinger liquid
\begin{align}
H_0 &= \frac{\vf}{4\pi} \intr{\sum_\ell \normord{ \kla*{\partial_x\varphi_\ell}^2 }}{x} = \frac{\vf}{2\pi} \intr{ \normord{ \klb*{\kla*{\partial_x\phi}^2 + \kla*{\partial_x\theta}^2} }}{x}
\end{align}
where we introduced the fields $\varphi_\ell\equiv \ell\phi - \theta$.
\subsection{Interaction}
We have already demonstrated that the terms $\psi_\ell^\dagger\psi_\ell$ become derivatives of the chiral fields.
There is no normal-ordering built into the interaction, but note that summing on $\ell$ also cancels the divergent constant since
\begin{align}
\sum_\ell \psi_\ell^\dagger(x) \psi_\ell(x) &= \lim_{\epsilon\rightarrow 0} \sum_\ell \klb*{\frac{-\i\ell}{2\pi\epsilon} - \frac{\ell}{2\pi}\partial_x\varphi_\ell(x)} = \frac{-1}{\pi}\partial_x\phi(x).
\end{align}
The other type of term appearing in $H_\Text{int}$ is of the form
\begin{align}
\psi_\ell^\dagger(x)\psi_{-\ell}(x) &= \frac{\eta_\ell \eta_{-\ell}}{L} \normord{\powe{\i\varphi_\ell(x)}} \normord{\powe{-\i\varphi_{-\ell}(x)}} = \frac{\eta_\ell \eta_{-\ell}}{L} \normord{\powe{2\i\ell\phi(x)}}
\end{align}
and is evidently much less difficult to bosonize due to the fact that $\varphi_\ell$ commutes for distinct species.
%In order to avoid mistakes with normal-ordering one should be careful not to naively insert these ``identities'' into more complicated expressions.
%To understand what can go wrong, consider the first Hamilton density\footnote{For brevity, limits that arise from point-splitting such as $\lim_{\epsilon\rightarrow 0}$ are implied}
%\begin{align}
%\mathcal{H}_{+,\ell,\ell'}^\Text{I} &= \psi_{\ell}^\dagger(x) \psi_{\ell}(x+\epsilon) \psi_{\ell'}^\dagger(x+a) \psi_{\ell'}(x+a+\epsilon)\\
%&= \frac{1}{L^2} \normord{\powe{\i\varphi_\ell(x)}} \normord{\powe{-\i\varphi_\ell(x+\epsilon)}} \normord{\powe{\i\varphi_{\ell'}(x+a)}} \normord{\powe{-\i\varphi_{\ell'}(x+a+\epsilon)}}\\
%&= \frac{1}{L^2} \normord{\powe{\i\varphi_\ell(x)-\i\varphi_\ell(x+\epsilon)}} \normord{\powe{\i\varphi_{\ell'}(x+a)-\i\varphi_{\ell'}(x+a+\epsilon)}} \powe{\klb*{\i\varphi_{\ell}^{(-)}(0), -\i \varphi_{\ell}^{(+)}(\epsilon)}} \powe{\klb*{\i \varphi_{\ell'}^{(-)}(0), -\i \varphi_{\ell'}^{(+)}(\epsilon)}}\\
%&= \frac{1}{2\pi\i\ell\epsilon} \frac{1}{2\pi\i\ell'\epsilon} \normord{\powe{\i\varphi_\ell(x)-\i\varphi_\ell(x+\epsilon) + \i\varphi_{\ell'}(x+a)-\i\varphi_{\ell'}(x+a+\epsilon)}} \powe{\klb*{\i\varphi_\ell^{(-)}(0)-\i\varphi_\ell^{(-)}(\epsilon), \i\varphi_{\ell'}^{(+)}(a)-\i\varphi_{\ell'}^{(+)}(a+\epsilon)}}\\
%&= \frac{-\ell\ell'}{4\pi^2\epsilon^2} \normord{\powe{\i\varphi_\ell(x)-\i\varphi_\ell(x+\epsilon) + \i\varphi_{\ell'}(x+a)-\i\varphi_{\ell'}(x+a+\epsilon)}}  \klb*{\frac{\i\ell a}{\i\ell (a-\epsilon)} \frac{\i\ell a}{\i\ell (a+\epsilon)}}^{\delta_{\ell,\ell'}}\\
%&= \frac{-\ell\ell'}{4\pi^2\epsilon^2} \normord{\powe{-\i\epsilon\partial_x\varphi_\ell(x)-\i\frac{\epsilon^2}{2} \partial_x^2\varphi_\ell(x) -\i\epsilon \partial_x \varphi_{\ell'}(x+a)-\i\frac{\epsilon^2}{2} \partial_x^2 \varphi_{\ell'}(x+a) }}
%\end{align}
%Since we integrate over the density we may drop all total derivatives.
%Taylor expanding then yields
%\begin{align}
%\mathcal{H}_{+,\ell,\ell'}^\Text{I} &= \frac{-\ell\ell'}{4\pi^2\epsilon^2} \normord{\klb*{1 - \frac{\epsilon^2}{2} \kla*{\partial_x\varphi_\ell(x)}^2 - \frac{\epsilon^2}{2} \kla*{\partial_x\varphi_{\ell'}(x+a)}^2 - \frac{\epsilon^2}{2} \partial_x\varphi_\ell(x) \partial_x \varphi_{\ell'}(x+a)}}.
%\end{align}
%Summing on $\ell$ and $\ell'$ removes the divergent constant and leaves us with
%\begin{align}
%\sum_{\ell,\ell'} \mathcal{H}_{+,\ell,\ell'}^\Text{I} &= \frac{1}{8\pi^2} \sum_{\ell,\ell'} \normord{\partial_x \phi(x) \partial_x\phi(x+a)} \simeq \frac{1}{2\pi^2} \normord{\kla*{\partial_x\phi}^2}.
%\end{align}
%The naive method looks correct up to a factor of two because
%\begin{align}
%\sum_{\ell,\ell'} \psi_{\ell}^\dagger(x) \psi_{\ell}(x+\epsilon) \psi_{\ell'}^\dagger(x+a) \psi_{\ell'}(x+a+\epsilon) &= \frac{-1}{\pi} \partial_x \phi(x) \frac{-1}{\pi} \partial_x\phi(x+a).
%\end{align}
%Both calculations are perfectly valid, but the naive result is not normal-ordered, and to avoid mistakes one should  generally take the safer route.\\
The first Hamiltonian density is rather straightforward to bosonize and yields
\begin{align}
\sum_{\ell,\ell'} \psi_{\ell}^\dagger(x) \psi_{\ell}(x+\epsilon) \psi_{\ell'}^\dagger(x+a) \psi_{\ell'}(x+a+\epsilon) &= \frac{-1}{\pi} \partial_x \phi(x) \frac{-1}{\pi} \partial_x\phi(x+a) \simeq \frac{1}{\pi^2} \normord{\kla*{\partial_x\phi}^2}.
\end{align}
In the last step we assumed that $a$ is small and dropped a constant by introducing the normal-ordering.
The second density becomes
\begin{align}
\mathcal{H}_{+,\ell,\ell'}^\Text{II} &= \psi_{\ell}^\dagger(x) \psi_{-\ell}(x) \psi_{\ell'}^\dagger(x+a) \psi_{-\ell'}(x+a)\\
&= \frac{\eta_\ell\eta_{-\ell} \eta_{\ell'}\eta_{-\ell'}}{L^2} \normord{\powe{\i\varphi_\ell(x) - \i\varphi_{-\ell}(x)}} \normord{\powe{\i\varphi_{\ell'}(x+a) - \i\varphi_{-\ell'}(x+a)}}\\
&= \frac{-\ell\ell'}{L^2} \normord{\powe{\i\varphi_\ell(x) - \i\varphi_{-\ell}(x) + \i\varphi_\ell(x+a) - \i\varphi_{-\ell}(x+a)}} \powe{\ell\ell' \klb*{\i \varphi_{\ell}^{(-)}(0), \i \varphi_{\ell}^{(+)}(a)}} \powe{\ell\ell' \klb*{-\i \varphi_{-\ell}^{(-)}(0), -\i \varphi_{-\ell}^{(+)}(a)}}\\
&= \frac{-\ell\ell'}{L^2} \klb*{\frac{2\pi\i\ell a}{L} \frac{-2\i\pi\ell a}{L}}^{\ell\ell'} \normord{\powe{2\i\ell\phi(x) + 2\i\ell'\phi(x+a)}}.
\end{align}
At this point we should consider two separate cases.
For $\ell=\ell'$ we find
\begin{align}
\mathcal{H}_{+,\ell,\ell}^\Text{II} &= \frac{-4\pi^2a^2}{L^4} \normord{\powe{2\i\ell\phi(x) + 2\i\ell'\phi(x+a)}} \simeq \frac{-4\pi^2a^2}{L^4} \normord{\powe{4\i\ell\phi(x)}}.
\end{align}
This term is negligible due to the prefactor of $\frac{1}{L^4}$.
For $\ell=-\ell'$ we instead find
\begin{align}
\mathcal{H}_{+,\ell,-\ell}^\Text{II} &= \frac{1}{4\pi^2a^2} \normord{\powe{2\i\ell\phi(x) - 2\i\ell\phi(x+a)}} \simeq \frac{1}{4\pi^2a^2} \normord{\powe{-2\i\ell a\partial_x\phi(x) - \i\ell a^2\partial_x^2\phi(x)}}\\
&\simeq \frac{1}{4\pi^2a^2} \normord{\klb*{1 - 2\i\ell a\partial_x\phi(x) - \i\ell a^2\partial_x^2\phi(x) - \frac{1}{2} (2\i a)^2\kla*{\partial_x\phi(x)}^2 }} \simeq \frac{1}{2\pi^2} \normord{\kla*{\partial_x\phi}^2}
\end{align}
up to constants and total derivatives.
%The second density should be considered in two separate cases.
%For $\ell=\ell'$ we find
%\begin{align}
%\mathcal{H}_{+,\ell,\ell}^\Text{II} &= \psi_{\ell}^\dagger(x) \psi_{-\ell}(x) \psi_{\ell}^\dagger(x+a) \psi_{-\ell}(x+a)\\
%&= \frac{\eta_\ell\eta_{-\ell} \eta_\ell\eta_{-\ell}}{L^2} \normord{\powe{\i\varphi_\ell(x) - \i\varphi_{-\ell}(x)}} \normord{\powe{\i\varphi_\ell(x+a) - \i\varphi_{-\ell}(x+a)}}\\
%&= \frac{-1}{L^2} \normord{\powe{\i\varphi_\ell(x) - \i\varphi_{-\ell}(x) + \i\varphi_\ell(x+a) - \i\varphi_{-\ell}(x+a)}} \powe{\klb*{\i \varphi_{\ell}^{(-)}(0), \i \varphi_{\ell}^{(+)}(a)}} \powe{\klb*{-\i \varphi_{-\ell}^{(-)}(0), -\i \varphi_{-\ell}^{(+)}(a)}}\\
%&= \frac{-1}{L^2} \frac{2\pi\i\ell a}{L} \frac{-2\i\pi\ell a}{L} \normord{\powe{2\i\ell\phi(x) + 2\i\ell\phi(x+a)}}\\
%&\simeq \frac{-4\pi^2a^2}{L^4} \normord{\powe{4\i\ell\phi(x)}}.
%\end{align}
%Due to the prefactor of $\frac{1}{L^4}$ this term is negligible.
%For $\ell=-\ell'$ we instead find
%\begin{align}
%\mathcal{H}_{+,\ell,-\ell}^\Text{II} &= \psi_{\ell}^\dagger(x) \psi_{-\ell}(x) \psi_{-\ell}^\dagger(x+a) \psi_{\ell}(x+a)\\
%&= \frac{\eta_\ell\eta_{-\ell} \eta_{-\ell}\eta_\ell}{L^2} \normord{\powe{\i\varphi_\ell(x) - \i\varphi_{-\ell}(x)}} \normord{\powe{\i\varphi_{-\ell}(x+a) - \i\varphi_{\ell}(x+a)}}\\
%&= \frac{1}{L^2} \normord{\powe{\i\varphi_\ell(x) - \i\varphi_{-\ell}(x) + \i\varphi_{-\ell}(x+a) - \i\varphi_{\ell}(x+a)}} \powe{\klb*{\i \varphi_{\ell}^{(-)}(0), -\i \varphi_{\ell}^{(+)}(a)}} \powe{\klb*{-\i \varphi_{-\ell}^{(-)}(0), \i \varphi_{-\ell}^{(+)}(a)}}\\
%&\simeq \frac{1}{2\pi\i\ell a} \frac{1}{-2\pi\i\ell a} \normord{\powe{-\i a\partial_x \varphi_\ell(x) - \i\frac{a^2}{2} \partial_x^2 \varphi_\ell(x) + \i a\partial_x \varphi_{-\ell}(x) + \i \frac{a^2}{2} \partial_x \varphi_{-\ell}(x)}}.
%\end{align}
%Up to a constant and total derivatives -- which we will once again ignore -- this can be simplified to
%\begin{align}
%\sum_\ell \mathcal{H}_{+,\ell,-\ell}^\Text{II} &\simeq \frac{1}{4\pi^2a^2} \sum_\ell \normord{\klb*{-\frac{a^2}{2} \kla*{\partial_x\varphi_\ell}^2 - \frac{a^2}{2} \kla*{\partial_x\varphi_{-\ell}}^2 + \frac{a^2}{2} \partial_x\varphi_\ell \partial_x \varphi_{-\ell}}}\\
%&= \frac{-1}{8\pi^2} \sum_\ell \normord{\klb*{ \kla*{\partial_x\phi}^2 + \kla*{\partial_x\theta}^2 - \kla*{\ell\partial_x\phi - \partial_x\theta} \kla*{-\ell\partial_x\phi - \partial_x\theta} }} = \frac{-1}{2\pi^2} \normord{\kla*{\partial_x\phi}^2}.
%\end{align}
The first $U_-$-part becomes
\begin{align}
\sum_{\ell,\ell'} \mathcal{H}_{-,\ell,\ell'} &= \sum_{\ell'} \frac{\eta_{\ell'} \eta_{-\ell'}}{L} \normord{\powe{2\i\ell'\phi(x+a)}}\frac{-1}{\pi} \partial_x\phi(x).
\end{align}
Under the integral we may now shift the argument of the fields by $a$.
Combining with the second part then yields
\begin{align}
H_- &= \frac{U_-}{2} \sum_\ell \frac{ -\eta_\ell \eta_{-\ell}}{2\pi^2\alpha} \intr{  \powe{2\i\ell\phi(x)} \klb*{\partial_x\phi(x-a) + \partial_x\phi(x+a)} }{x}.
\end{align}
\subsection{Final model}
Except for neglecting the sine-Gordon term $\sim\cos(4\phi)$ we arrive at a low-energy theory described by $H=H_0 + gH_\Text{int}$ where
\begin{align}
H_0 &= \frac{u}{2\pi} \intr{\klb*{K \kla*{\partial_x\theta}^2 + \frac{1}{K} \kla*{\partial_x\phi}^2}}{x},\\
H_\Text{int} &= g\frac{u}{2\pi\alpha} \sum_\ell \eta_\ell \eta_{-\ell} \intr{\powe{2\i\ell\phi(x)} \klb*{\partial_x\phi(x+a) + \partial_x\phi(x-a)}}{x}
\end{align}
with the dimensionless coupling constant $g = \frac{-U_-a}{2\pi u}$, the renormalized Fermi velocity $u = \frac{\vf}{K}$ and the Luttinger parameter $K = \kla*{1 + \frac{2U_+ a}{\pi\vf}}^{-1/2}$.
\section{Green's function}
Before we dive into the calculations let us introduce the notation
\begin{align}
\bra*{\zabx{A_1&A_2&A_3&\dots}{B_1&B_2&B_3&\dots}{z_1&z_2&z_3&\dots}}_0 &= \bra[\bigg]{\prod_j \powe{\i A_j\phi(z_j) + \i B_j\theta(z_j)}}_0.
\end{align}
We can rewrite this expression as
\begin{align}
\bra*{\abx{A}{B}{Z}}_0 &= \prod_{j<k} \powe{\frac{1}{2}\klb*{KA_jA_k + \frac{1}{K}B_jB_k + \abs*{s_{jk}}}F_1(z_j-z_k)} \powe{-\frac{1}{2}|s_{jk}| F\kla*{z_{-\sgn(s_{jk}),j} - z_{-\sgn(s_{jk}), k}}}
\end{align}
where $s_{jk} = A_jB_k+A_kB_j$ and the two functions are
\begin{align}
F(z_\ell) &= \log\klb*{\frac{\beta u}{\pi\alpha}} + \log\klb*{\sin(z_\ell)},\\
F_1(z) &= \frac{1}{2} \klb*{F(z_+) + F(z_-)}
\end{align}
with the coordinates $z_\pm = \frac{\pi\tau}{\beta} \pm \i \frac{\pi x}{\beta u}$ and the tuple $z=(z_+,z_-)$.
\subsection{Zeroth order}
The Matsubara Green's function is given by
\begin{align}
G_{0,\ell,\ell'}^\M(z) &= -\mathcal{T} \bra*{\psi_\ell(z) \psi_{\ell'}^\dagger(0)}_0 = -\frac{1}{2\pi a} \bra*{\zabx{-\ell & \ell'}{1 & -1}{z & 0}}_0\\
&= \frac{-\delta_{\ell,\ell'}}{2\pi a} \powe{\frac{1}{2}\klb*{-K - \frac{1}{K} + 2}F_1(z)} \powe{-F\kla*{z_{-\ell}}} = \frac{-\delta_{\ell,\ell'}}{2\beta u} \kla*{\frac{\pi\alpha}{\beta u}}^M \csc(z_{-\ell})^{1 + \frac{M}{2}} \csc(z_\ell)^\frac{M}{2}
\end{align}
where $M=\frac{1}{2} \klb*{K + \frac{1}{K} - 2} \ge 0$.
At this point we give the Fourier integrals
\begin{align}
J_0(k,\i\omega_n^\F; b) &= \intx{0}{\beta}{\intr{\powe{\i \omega_n^\F\tau - \i kx} \csc(z_+)^{1+b} \csc(z_-)^b}{x} }{\tau}\\
&= -\i\sin(\pi b) \frac{u\beta^2}{\pi^2} I_0\kla*{\beta\frac{\i\omega_n^\F + uk}{2\pi}, b} I_0\kla*{\beta\frac{\i\omega_n^\F - uk}{2\pi}, b+1},\\
L_0(k,\i\omega_n^\B; b) &= \intx{0}{\beta}{\intr{\powe{\i \omega_n^\B \tau - \i kx} \csc(z_+)^{b} \csc(z_-)^b}{x} }{\tau}\\
&= \sin(\pi b) \frac{u\beta^2}{\pi^2} I_0\kla*{\beta\frac{\i\omega_n^\B + uk}{2\pi}, b} I_0\kla*{\beta\frac{\i\omega_n^\B - uk}{2\pi}, b}
\end{align}
where $I_0(z,b) = \frac{2^{b-1}}{\pi} \Gamma(1-b) \Gamma\kla*{\frac{b+\i z}{2}} \Gamma\kla*{\frac{b-\i z}{2}} \sin\kla*{\frac{\pi b}{2} + \frac{\pi\i z}{2}}$.
%#beta csc b+1,b
%def i_0(z, b):
%    return (2**(b-1) * special.gamma(b/2+1j*z/2) * special.gamma(b/2-1j*z/2)* np.sin(np.pi*b/2 + 1j*z*np.pi/2) * special.gamma(1-b) / np.pi)
%def j_0(k, i_omega_n, beta=1.1, b=0.5, u=1.0):
%    return (-1j * u * beta**2 / np.pi**2 * np.sin(np.pi*b) * i_0(beta * (i_omega_n - u*k) / (2*np.pi), b+1) * i_0(beta * (i_omega_n + u*k) / (2*np.pi), b) )
%def l_0(k, i_omega_n, beta=1.1, b=0.5, u=1.0):
%    return (u * beta**2 / np.pi**2 * np.sin(np.pi*b) * i_0(beta * (i_omega_n - u*k) / (2*np.pi), b) * i_0(beta * (i_omega_n + u*k) / (2*np.pi), b) )
%beta=2.4; u=1.3
%n=2; om = 1j*(2*n+1)*np.pi/beta; k=0.83; b=0.1; L=10
%print(c_quad(lambda tau: np.exp(om*tau) * c_quad(lambda x: np.exp(-1j*k*x) / np.sin(np.pi/(beta*u)*(tau*u+1j*x))**(b+1) / np.sin(np.pi/(beta*u)*(tau*u-1j*x))**b, -L, L)[0], 0, beta))
%print(j_0(k,om,beta,b,u))
%om = 1j*(2*n)*np.pi/beta
%print(c_quad(lambda tau: np.exp(om*tau) * c_quad(lambda x: np.exp(-1j*k*x) / np.sin(np.pi/(beta*u)*(tau*u+1j*x))**b / np.sin(np.pi/(beta*u)*(tau*u-1j*x))**b, -L, L)[0], 0, beta))
%print(l_0(k,om,beta,b,u))
Note that up to prefactors $J_0$ already corresponds to the Fourier transform of the free Green's function.
\subsection{First order}
Using the replacement
\begin{align}
\partial_x\phi(x+a) &= \lim_{\epsilon\rightarrow a} \partial_\epsilon \lim_{J\rightarrow 0} (-\i\partial_J) \powe{\i J\phi(x+\epsilon)}
\end{align}
allows us to evaluate the first part of the interaction to\footnote{$\tilde{z}' = z' + \epsilon$}
\begin{align}
G_{1,\ell,\ell'}^\M(z) &= \frac{1}{(2\pi a)^2} \lim_{\epsilon\rightarrow a} \partial_\epsilon (-\i\partial_J) \int \sum_s \eta_\ell \eta_s \eta_{-s} \eta_{\ell'} \bra*{\powe{-\i\varphi_\ell(z)} \powe{2\i s\phi(z')} \powe{\i J\phi(z'+\epsilon)} \powe{\i\varphi_{\ell'}(0)}}_0\,\d z'\\
&= \frac{\delta_{\ell,-\ell'}}{(2\pi a)^2} \lim_{\epsilon\rightarrow a} \partial_\epsilon (-\i\partial_J) \int \bra*{\zabx{-\ell & 2\ell & J & -\ell}{1 & 0 & 0 & -1}{z & z' & \tilde{z}' & 0}}_0\,\d z'\\
&= \frac{\delta_{\ell,-\ell'}}{(2\pi a)^2} \lim_{\epsilon\rightarrow a} \partial_\epsilon (-\i\partial_J) \bsplit{&\powe{(1-K) F_1(z-z') -F\kla[\big]{z_{-\ell} - z_{-\ell}'}} \powe{(1-K) F_1(z-z')-F\kla*{z_{\ell}'}} \powe{-NF_1(z)} \powe{-\ell JF_1(\epsilon)}\\
&\times \powe{\frac{1}{2}\klb*{-\ell JK + J}F_1(z-\tilde{z}')} \powe{-\frac{1}{2}J F\kla*{z_{-} - \tilde{z}_{-}'}} \powe{\frac{1}{2}\klb*{-\ell JK + J}F_1(\tilde{z}')} \powe{-\frac{1}{2}J F\kla*{\tilde{z}_{+}'}}}\\
&= \frac{\ell\i \delta_{\ell,-\ell'}}{2(2\pi a)^2} \int \bsplit{&\powe{(1-K) F_1(z-z') -F\kla[\big]{z_{-\ell} - z_{-\ell}'}} \powe{(1-K) F_1(z-z')-F\kla*{z_{\ell}'}} \powe{-NF_1(z)}\\
&\times \partial_{\tilde{x}'} \klb*{\kla*{ K-\ell}F_1(z-\tilde{z}')+\ell F\kla*{z_{-} - \tilde{z}_{-}'}+\kla*{ K -\ell}F_1(\tilde{z}')+\ell F\kla*{\tilde{z}_{+}'}}\,\d z'}
\end{align}
where $N=\frac{1}{2} \klb*{\frac{1}{K} - K}$.
The time-ordering would in principle require us to split the $\tau'$-integration into an interval from $0$ to $\tau$ and one from $\tau$ to $\beta$.
On the second interval the integrand would pick up an additional minus sign from the reordered Klein factors and we have to replace $z-z'$ by $z'-z$.
Since the latter amounts to another minus sign, the integrand remains the same.
Lastly, the term containing the derivatives changes from $\sim F(z_\ell - \tilde{z}_\ell')$ to $\sim F(\tilde{z}_\ell' - z_\ell)$.
From the definition of the functions we see that corresponds to adding $\pm\pi\i$, and this constant vanishes upon differentiation.
Thus we can recombine simply recombine the two integrations into the one given above.\\
In order to obtain the full correction we now need to replace $\epsilon=\pm a$ and add up both terms.
This is also why we ignored the term $F'(\epsilon)$ that drops out in this process due to the anti-symmetry of $F'$.
Note that the term containing the derivatives would lead to a total derivative upon setting $a=0$.
To leading order we should therefore expand in $a$ giving
\begin{align}
G_{1,\ell,\ell'}^\M(z) &= \ell\i \frac{\delta_{\ell,-\ell'}}{(2\pi)^2} \int \bsplit{&\powe{(1-K) F_1(z-z') -F\kla[\big]{z_{-\ell} - z_{-\ell}'}} \powe{(1-K) F_1(z-z')-F\kla*{z_{\ell}'}} \powe{-NF_1(z)}\\
&\times \partial_{x'}^3 \klb*{\kla*{ K-\ell}F_1(z-z')+\ell F\kla*{z_{-} - z_{-}'}+\kla*{ K -\ell}F_1(z')+\ell F\kla*{z_{+}'}}\,\d z'}
\end{align}
The term containing the derivatives may be rewritten as
\begin{align}
\sum_{s=\pm 1} &\partial_{x'}^3 \klb*{ \kla*{\frac{K+s\ell}{2}} \log\klb*{\sin\kla*{z_{-s} - z_{-s}'}} + \kla*{\frac{K+s\ell}{2}} \log\klb*{\sin\kla*{z_s'}}}\\
&=\sum_{s=\pm 1} \frac{K+s\ell}{2} \partial_{x'}^2 \klb*{ \frac{\pi \i s}{\beta u} \cot\kla*{z_{-s} - z_{-s}'} + \frac{\pi \i s}{\beta u} \cot\kla*{z_s'}}\\
&= -\sum_{s=\pm 1} \i \ell s \kla*{K+s} \frac{\pi^3}{\beta^3 u^3} \klb*{ \cot\kla*{z_{-\ell s} - z_{-\ell s}'} \csc\kla*{z_{-\ell s} - z_{-\ell s}'}^2 + \cot\kla*{z_{\ell s}'} \csc\kla*{z_{\ell s}'}^2}.
\end{align}
The full first order Green's function thus becomes
\begin{align}
G_{1,\ell,-\ell}^\M(z) &= C_1\powe{-NF_1(z)} \int \bsplit{& \sum_{s=\pm 1} \kla*{sK+1} \csc\kla*{z_{-\ell}'}^{\frac{K-1}{2}} \csc\kla*{z_{\ell}'}^{\frac{K+1}{2}}\\
&\times \csc\kla*{z_{-\ell} - z_{-\ell}'}^{\frac{K+1}{2}} \csc\kla*{z_{\ell} - z_{\ell}'}^{\frac{K-1}{2}} \\
&\times \klb[\Big]{ \cot\kla*{z_{-\ell s} - z_{-\ell s}'} \csc\kla*{z_{-\ell s} - z_{-\ell s}'}^2 + \cot\kla*{z_{\ell s}'} \csc\kla*{z_{\ell s}'}^2}\,\d z'.}
\end{align}
where $C_1 = \frac{1}{4\pi^2} \kla*{\frac{\pi a}{\beta u}}^{N+2K} \frac{\pi^3}{\beta^3 u^3}$.
This result is somewhat easier to obtain by expanding the field within the interaction Hamiltonian itself.
The disadvantage of this approach is that it is less apparent why the term $F'(a)$ should vanish.\\
The Fourier transform can be simplified by inserting a resolution of unity via
\begin{align}
G_{1,\ell,-\ell}^\M &= C_1 \int \bsplit{& \powe{\i (2n+1)(\tau+\tau') - \i k(x+x')} \powe{-NF_1(z'')} \csc\kla*{z_{-\ell}}^{\frac{K+1}{2}} \csc\kla*{z_{\ell}}^{\frac{K-1}{2}} \\
&\times \csc\kla*{z_{-\ell}'}^{\frac{K-1}{2}} \csc\kla*{z_{\ell}'}^{\frac{K+1}{2}} \delta(\tau''-\tau-\tau') \delta(x''-x-x')\\
&\times \sum_{s=\pm 1} \kla*{sK+1} \klb[\Big]{ \cot\kla*{z_{-\ell s}}\csc\kla*{z_{-\ell s}}^2 + \cot\kla*{z_{\ell s}'} \csc\kla*{z_{\ell s}'}^2 }\,\d z'' \,\d z'\,\d z'}\\
&= \frac{C_1}{2\pi\beta} \int \bsplit{& \sum_{s=\pm 1} \sum_{m\in\mathbb{Z}} \kla*{sK+1} \powe{\i \omega_{m}^\B \tau'' -\i k'x''} \powe{-NF_1(z'')}\\
&\times \klb[\bigg]{\powe{\i \omega_{n-m}^\F \tau' -\i (k-k') x'}  \csc\kla*{z_{-\ell}'}^{\frac{K-1}{2}} \csc\kla*{z_{\ell}'}^{\frac{K+1}{2}}\\
&\qquad\times \powe{\i \omega_{n-m}^\F \tau -\i (k-k') x} \csc\kla*{z_{-\ell}}^{\frac{K+1}{2}} \csc\kla*{z_{\ell}}^{\frac{K-1}{2}} \cot\kla*{z_{-\ell s}} \csc\kla*{z_{-\ell s}}^2\\
&\quad + \powe{\i \omega_{n-m}^\F \tau' -\i (k-k') x'}  \csc\kla*{z_{-\ell}'}^{\frac{K-1}{2}} \csc\kla*{z_{\ell}'}^{\frac{K+1}{2}}\\
&\quad\quad\times \powe{\i \omega_{n-m}^\F \tau -\i (k-k') x} \csc\kla*{z_{-\ell}}^{\frac{K+1}{2}} \csc\kla*{z_{\ell}}^{\frac{K-1}{2}} \cot\kla*{z_{\ell s}'} \csc\kla*{z_{\ell s}'}^2} \,\d (z,z',z'')\,\d k' }\\
&= \frac{C_1}{2\pi\beta} \bsplit{\intr{& \sum_{s=\pm 1} \sum_{m\in\mathbb{Z}} (sK+1) L_0\kla*{k',\i \omega_m^\B; \frac{N}{2}} \\
&\times \klb[\bigg]{ J_0\kla*{\ell(k-k'), \i \omega_{n-m}^\F; \frac{K-1}{2}}  J_1\kla*{-\ell(k-k'), \i \omega_{n-m}^\F; \frac{K-1}{2}, s}\\
&\qquad + J_1\kla*{\ell(k-k'), \i \omega_{n-m}^\F; \frac{K-1}{2}, s} J_0\kla*{-\ell(k-k'), \i \omega_{n-m}^\F; \frac{K-1}{2}}} }{k'}}
\end{align}
where we introduced the functions
\begin{align}
J_1(k,\i\omega_n^\F;b,s) &= \intx{0}{\beta}{\intr{\powe{\i\omega_n^\F\tau - \i kx} \csc\kla*{z_+}^{1+b} \csc\kla*{z_-}^b \csc\kla*{z_s}^2 \cot\kla*{z_s} }{x} }{\tau}.
\end{align}
The problem is thus reduced to that of finding an -- exact or at least approximate -- analytic expression for $J_1$.
It is best to rewrite the function as $J_{1,s}$ where
\begin{align}
J_{1,+}(k,\i\omega_n^\F; b) &= \intx{0}{\beta}{\intr{\powe{\i \omega_n^\F\tau - \i kx} \csc\kla*{z_+}^{3+b} \csc\kla*{z_-}^{b} \cot\kla*{z_+} }{x} }{\tau}\\
&= -\frac{1}{3+b} \frac{\beta}{\pi} \intx{0}{\beta}{\intr{\powe{\i \omega_n^\F\tau - \i kx} \csc\kla*{z_-}^{b} \partial_\tau \csc\kla*{z_+}^{3+b} }{x} }{\tau},\\
J_{1,-}(k,\i\omega_n^\F; b) &= \intx{0}{\beta}{\intr{\powe{\i \omega_n^\F\tau - \i kx} \csc\kla*{z_+}^{1+b} \csc\kla*{z_-}^{2+b} \cot\kla*{z_-} }{x} }{\tau}\\
&= -\frac{1}{2+b} \frac{\beta}{\pi} \intx{0}{\beta}{\intr{\powe{\i \omega_n^\F\tau - \i kx} \csc\kla*{z_+}^{1+b} \partial_\tau \csc\kla*{z_-}^{2+b} }{x} }{\tau}.
\end{align}
Expressed in terms of these functions we now find
\begin{align}
G_{1,\ell,-\ell}^\M(k,\i\omega_n^\F) &= \frac{C_1}{2\pi\beta} \sum_{s,s'} (sK+1) \bsplit{\intr{ &\sum_{m\in\mathbb{Z}}  L_0\kla*{k', \i \omega_m^\B; \frac{N}{2}} J_0\kla*{s'(k-k'), \i \omega_{n-m}^\F; \frac{K-1}{2}}\\
&\times J_{1,s}\kla*{-s'(k-k'), \i \omega_{n-m}^\F; \frac{K-1}{2}} }{k'}.}
\end{align}
%%It is best to rewrite the function as $J_{1,s}$ where
%%\begin{align}
%%J_{1,+}(k,\i\omega_n^\F; b) &= \intx{0}{\pi}{\intr{\powe{\i (2n+1)\tau - \i kx} \csc\kla*{\tau + \i x}^{1+b} \csc\kla*{\tau - \i x}^{b} \cot\kla*{\tau + \i x} }{x} }{\tau}\\
%%&= -\frac{1}{1+b} \intx{0}{\pi}{\intr{\powe{\i (2n+1)\tau - \i kx} \csc\kla*{\tau - \i x}^{b} \partial_\tau \csc\kla*{\tau + \i x}^{1+b} }{x} }{\tau},\\
%%J_{1,-}(k,\i\omega_n^\F; b) &= \intx{0}{\pi}{\intr{\powe{\i (2n+1)\tau - \i kx} \csc\kla*{\tau + \i x}^{1+b} \csc\kla*{\tau - \i x}^{b} \cot\kla*{\tau - \i x} }{x} }{\tau}\\
%%&= -\frac{1}{b} \intx{0}{\pi}{\intr{\powe{\i (2n+1)\tau - \i kx} \csc\kla*{\tau + \i x}^{1+b} \partial_\tau \csc\kla*{\tau - \i x}^{b} }{x} }{\tau}.
%%\end{align}
%%Partial integration on $\tau$ and rewriting the derivative $\partial_\tau$ as $\pm\i\partial_x$ followed by partial integration on $x$ allows us to derive the two relations
%%\begin{align}
%%(1+b)J_{1,+} + bJ_{1,-} &= \i\omega_n^\F J_0,\qquad (1+b)J_{1,+} - bJ_{1,-} = -k J_0.
%%\end{align}
%%Rearranging then yields the exact solutions $J_{1,+} = \frac{\i\omega_n^\F - k}{2(1+b)} J_0$ and $J_{1,-} = \frac{\i\omega_n^\F + k}{2b} J_0$.
%%%n=2; om = 1j*(2*n+1); k=0.43; K=0.9; b=(K-1)/2; L=10
%%%print(c_quad(lambda tau: np.exp(1j*(2*n+1)*tau) * c_quad(lambda x: np.exp(-1j*k*x) / np.sin(tau+1j*x)**(1+b) / np.sin(tau-1j*x)**b/np.tan(tau-1j*x), -L, L)[0], 0, np.pi))
%%%print(c_quad(lambda tau: np.exp(1j*(2*n+1)*tau) * c_quad(lambda x: np.exp(-1j*k*x) / np.sin(tau+1j*x)**(1+b) / np.sin(tau-1j*x)**b/np.tan(tau+1j*x), -L, L)[0], 0, np.pi))
%%%j_0(k,om,np.pi,b) * (om-k)/2 / (1+b), j_0(k,om,np.pi,b) * (om+k)/2 / (b)
Exactly evaluating $J_{1,-}$ is surprisingly simple.
Integration by parts and replacing $\partial_\tau$ by $\pm\i\partial_x$ gives
\begin{align}
J_{1,-} &= \frac{\i\omega_n^\F}{2+b} \frac{\beta}{\pi} J_0\kla*{-k,\i\omega_n^\F; 1+b} + \frac{1}{2+b} \frac{\beta}{\pi} \intx{0}{\beta}{\intr{\powe{\i\omega_n^\F\tau - \i kx} \csc(z_-)^{2+b} \partial_\tau \csc(z_+)^{1+b}}{x} }{\tau}\\
&= \frac{\i\omega_n^\F}{2+b} \frac{\beta}{\pi} J_0\kla*{-k,\i\omega_n^\F; 1+b} + \frac{-\i u}{2+b} \frac{\beta}{\pi} \intx{0}{\beta}{\intr{\powe{\i\omega_n^\F\tau - \i kx} \csc(z_-)^{2+b} \partial_x \csc(z_+)^{1+b}}{x} }{\tau}\\
&= \frac{\i\omega_n^\F + uk}{2+b} \frac{\beta}{\pi} J_0\kla*{-k,\i\omega_n^\F; 1+b} + \frac{\i u}{2+b} \frac{\beta}{\pi} \intx{0}{\beta}{\intr{\powe{\i\omega_n^\F\tau - \i kx} \csc(z_+)^{1+b} \partial_x \csc(z_-)^{2+b}}{x} }{\tau}.
\end{align}
The remaining integral is once again $J_{1,-}$ which then yields $J_{1,-}(k,\i\omega_n^\F;b) = \frac{\i\omega_n^\F + uk}{2(2+b)} \frac{\beta}{\pi} J_0(-k,\i\omega_n^\F; 1+b)$.
%Via similar manipulations we may show that
%\begin{align}
%J_{1,+} &= \frac{\i\omega_n^\F - k}{3+b} \intx{0}{\pi}{\intr{\powe{\i\omega_n^\F \tau - \i kx} \csc(\tau - \i x)^b \csc(\tau + \i x)^{3+b}}{x} }{\tau}.
%\end{align}
%We may once again make use of our knowledge of $J_0$ to solve this by rewriting as
%\begin{align}
%J_{1,+} &= \frac{\i\omega_n^\F - k}{3+b} \intx{0}{\pi}{\intr{\powe{\i\omega_n^\F \tau - \i kx} \sin(\tau - \i x)^2 \csc(\tau - \i x)^{2+b} \csc(\tau + \i x)^{3+b}}{x} }{\tau}\\ 
%&= \frac{\i\omega_n^\F - k}{3+b} \intx{0}{\pi}{\intr{\powe{\i\omega_n^\F \tau - \i kx} \frac{2 - \powe{2\i\tau + 2x} - \powe{-2\i\tau - 2x}}{4} \csc(\tau - \i x)^{2+b} \csc(\tau + \i x)^{3+b}}{x} }{\tau}\\
%&=  \frac{\i\omega_n^\F - k}{3+b} \frac{1}{4}\klb*{2J_0\kla*{k,\i\omega_n^\F; 2+b} - J_0\kla*{k-2, \i\omega_n^\F + 2\i; 2+b} - J_0\kla*{k+2, \i\omega_n^\F - 2\i; 2+b} }.
%\end{align}
Computing $J_{1,+}$ is significantly more challenging.
However, assuming that the exponent $b$ is small, we may neglect the corresponding cosecant factor.
This restricts our approximation to $K$ being somewhat close to unity, but within this limit the integral becomes quite doable.
As shown before it gives
\begin{align}
J_{1,+}(k,\i\omega_n^\F;b) &\approx -\frac{1}{3+b}  \frac{\beta}{\pi} \intx{0}{\beta}{\intr{\powe{\i \omega_n^\F \tau - \i kx} \partial_\tau \csc(z_+)^{3+b}}{x} }{\tau}\\
%&= \frac{k}{3+b} \frac{\beta}{\pi} \intx{0}{\beta}{\powe{\i \omega_n^\F \tau} \frac{2^{b+3}}{\Gamma(b+3)} \abs*{\Gamma\kla*{\frac{b+3}{2}+\i \frac{\beta uk}{2\pi}}}^2 \powe{k\kla*{\tau - \frac{\pi}{2}}}  }{\tau}\\
&= \frac{uk}{3+b} \frac{u\beta^2}{\pi^2} \frac{2^{b+3}}{\Gamma(b+3)} \abs*{\Gamma\kla*{\frac{b+3}{2}+\i \frac{\beta uk}{2\pi}}}^2 \frac{\cosh\kla*{\frac{\beta u k}{2}}}{\i\omega_n^\F + uk}.
\end{align}
%#J1+
%n=2; u=1.0; beta=2.4; om = 1j*(2*n+1)*np.pi/beta; k=0.43; K=0.9; b=(K-1)/2; L=10; alpha=1e-3
%print(c_quad(lambda tau: np.exp(om*tau) * c_quad(lambda x: np.exp(-1j*k*x) / np.sin(np.pi/beta/u * (u*tau+1j*x))**(3+b) / np.tan(np.pi/beta/u * (u*tau+1j*x)), -L, L)[0], 0, beta))
%print(beta**2/np.pi**2 * u*k * 2**(b+3)/special.gamma(b+4) * np.cosh(beta*u*k/2) / (om + u*k) * np.abs(special.gamma((b+3+1j*beta*u*k/np.pi)/2))**2)
%\begin{align}
%J_{1,+}(-k) &= \intx{0}{\pi}{\intr{\powe{\i \omega_n^\F\tau - \i kx} \csc\kla*{\tau - \i x}^{3+b} \csc\kla*{\tau + \i x}^{b} \cot\kla*{\tau - \i x} }{x} }{\tau}\\
%&= \intx{0}{\pi}{\intr{\powe{\i \omega_n^\F\tau - \i kx} \csc\kla*{\tau - \i x}^{3+b} \csc\kla*{\tau + \i x}^{2+b} \cot\kla*{\tau - \i x} \sin\kla*{\tau + \i x}^2 }{x} }{\tau}\\
%&= \intx{0}{\pi}{\intr{\powe{\i \omega_n^\F\tau - \i kx} \csc\kla*{\tau - \i x}^{3+b} \csc\kla*{\tau + \i x}^{2+b} \cot\kla*{\tau - \i x} \frac{2 - \powe{2\i\tau - 2x} - \powe{-2\i\tau + 2x}}{4} }{x} }{\tau}\\
%&= \frac{1}{4} \klb*{2J_{1,-}\kla*{k,\i\omega_n^\F; 1+b} - J_{1,-}\kla*{k+2,\i\omega_n^\F+2\i; 1+b} - J_{1,-}\kla*{k-2,\i\omega_n^\F-2\i; 1+b}}.
%\end{align}
\section{Effective Hamiltonian}
The effective Hamiltonian is given by
\begin{align}
H_\Text{eff}(k,\omega) &= \omega\sigma_0 - \kla*{G^\R(k,\omega)}^{-1}
\end{align}
where $\sigma_0$ is the identity matrix.
Since $H_\Text{eff}\in\mathbb{C}^{2\times 2}$ we can decompose it into 
\begin{align}
H_\Text{eff}(k,\omega) &= d_0(k,\omega) \sigma_0 + \textbf{d}(k,\omega) \cdot \bm{\sigma}
\end{align}
where $\bm{\sigma} = (\sigma_x, \sigma_y, \sigma_z)^T$ and $\textbf{d} = (d_x,d_y,d_z)^T$.
The eigenvalues are given by
\begin{align}
E_\pm(k,\omega) &= d_0 \pm \sqrt{\textbf{d}\cdot \textbf{d}}.
\end{align}
One should be careful not to simplify the square-root to $|\textbf{d}|$ since the components are generally complex in our setting.
Writing $\textbf{d}=\textbf{d}_\r + \i \textbf{d}_\i$ instead yields
\begin{align}
E_\pm(k,\omega) &= d_0 \pm \sqrt{\textbf{d}_\r^2 - \textbf{d}_\i^2 + 2\i \textbf{d}_\r \cdot \textbf{d}_\i}.
\end{align}
Exceptional points are found at non-trivial solutions to the two equations
\begin{align}
\textbf{d}_\r^2 - \textbf{d}_\i^2 &= 0, \label{eqn:dr2 minus di2}\\
\textbf{d}_\r \cdot \textbf{d}_\i &= 0. \label{eqn:dr cdot di}
\end{align}
The numerical solution can be visualized as contour plots, see \autoref{fig:ep_order1} for an example.
\begin{figure}[!ht]
\centering
%\includegraphics{figures/ep_density_beta0.1_a0.01_vf0.2_Ua100.0_Ub-95.0_order1.pdf}
%\includegraphics{figures/ep_density_beta0.1_a0.05_vf0.2_Ua2.0_Ub-1.0_order1.pdf}
%\includegraphics{figures/ep_density_beta0.05_a0.005_vf0.05_Ua1.9_Ub-1.1_order1.pdf}
\includegraphics{figures/ep_density_beta0.05_a0.005_vf0.05_Ua1.5_Ub-1.3_order1.pdf}
\caption{Contour plot of \autoref{eqn:dr2 minus di2} and \autoref{eqn:dr cdot di} to first order with $a=0.005$, $\vf=0.05$, $\beta=0.05$, $U_\A=1.5$ and $U_\B=-1.3$.
This leads to $g\approx -0.022$ and $K\approx 0.997$, well within the realm of parameters where we can expect decent results.}
\label{fig:ep_order1}
\end{figure}
For the same set of parameters we may visualize the complex eigenvalues $E_\pm$.
In the static limit $\omega=0$ we find a pair of exceptional points for small momenta $k$ as is visualized in \autoref{fig:complex_E_omega0_order1}.
\begin{figure}[!ht]
\centering
%\includegraphics{figures/complex_E_density_beta0.1_a0.01_vf0.2_Ua100.0_Ub-95.0_omega0_order1.pdf}
%\includegraphics{figures/complex_E_density_beta0.1_a0.05_vf0.2_Ua2.0_Ub-1.0_omega0_order1.pdf}
%\includegraphics{figures/complex_E_density_beta0.05_a0.005_vf0.05_Ua1.9_Ub-1.1_omega0_order1.pdf}
\includegraphics{figures/complex_E_density_beta0.05_a0.005_vf0.05_Ua1.5_Ub-1.3_omega0_order1.pdf}
\caption{Complex energy eigenvalues $E_\pm$ of the effective Hamiltonian $H_\Text{eff}$ to first order for $a=0.005$, $\vf=0.05$, $\beta=0.05$, $U_\A=1.5$ and $U_\B=-1.3$.
The frequency is fixed to $\omega=0$ corresponding to the exceptional points in \autoref{fig:ep_order1}.}
\label{fig:complex_E_omega0_order1}
\end{figure}
%%% PLOT CODE ::
%beta=0.1; vf=0.2; ua=2.0; ub=-1; u_minus=(ua-ub)/2; u_plus=(ua+ub)/2
%alpha=5e-2; kwm = 0.005; wm=0.001
%K = 1/np.sqrt(1 + u_plus*alpha/(2*np.pi*vf)); u=vf/K; g=-u_minus*alpha/(2*np.pi*u)
%k_vals = np.linspace(-kwm, kwm, 613)
%omega_vals = np.array([0.0])
%green = green_perturbative(k_vals, omega_vals, beta, K, u, g, order=1, model="density var", alpha=alpha)
%hamilton = green_toolkit.hamilton_from_green(omega_vals, green)
%energies = np.array([green_toolkit.eigenvalues_2x2(hamilton[i,0]) 
%                     for i in range(k_vals.size)])
%energies = green_toolkit.convert_to_contiguous_arrays(energies.T[0], energies.T[1])
%fig, ax = plt.subplots()
%plot_complex(ax, k_vals, energies)
%ax.set_xlabel(r"$k$")
%ax.set_ylabel(r"$E_\pm$")
%ax.set_xlim(k_vals[0], k_vals[-1])
%mpl_special.embed_labels(fig, ax)
%print(energies[0].imag.max() - energies[1].imag.min(), energies[0][0], g, K)
%k_vals = np.linspace(-kwm, kwm, 201)
%omega_vals = np.linspace(-wm, wm, 180)
%green = green_perturbative(k_vals, omega_vals, beta, K, u, g, order=1, model="density var", alpha=alpha)
%hamilton = green_toolkit.hamilton_from_green(omega_vals, green)
%fig, ax = plt.subplots()
%ax.set_xlabel(r"$k$")
%ax.set_ylabel(r"$\omega$")
%ep_contour_plot(ax, hamilton, k_vals, omega_vals)
%eqn1, eqn2 = green_toolkit.hamilton_2x2_ep_eqn(hamilton)
%phase = np.sqrt(eqn1 + 2j*eqn2)
%arg = np.arctan2(phase.imag, phase.real)
%img = ax.imshow(arg.T, cmap="RdBu", aspect="auto", extent=[-kwm, kwm, -wm, wm], origin="lower", vmin=-np.pi/2, vmax=np.pi/2, alpha=0.85)
%cbar = fig.colorbar(img, ax=ax, label=r"$\arg(\Delta E)$")
%mpl_special.format_ticklabels(cbar.ax, which="y")
%mpl_special.embed_labels(fig, [ax, cbar.ax])

%\begin{thebibliography}{9}
%\bibitem{example} description, \url{https://www.exmapleurl}, day.\,month~2021.
%\end{thebibliography}

\end{document}